\section{结论}
\label{sec:conclusion}

本文描述 PBM：一个作用在 sRGB 显示值空间、参考校准的解析混色模型。模型将 K--M/Duncan 结构的核心与三层连续混合的修复机制结合：针对深色加白的逐通道幂律目标、针对互补色对的加性谐波修正，以及使模型交换对称的软角色分配。谐波修正是十七个比值特征上的线性模型，因此整个模型不含三角函数与对数，只保留算术、幂与平方根。它 $O(1)$、无查找表，在 $\alpha=0$ 与 $\alpha=1$ 处精确返回输入，按构造连续，并在浮点舍入范围内交换对称（实测探针 $6.7\times10^{-16}$）。51 个谐波系数由最小二乘闭式求得，其余 14 个标量由同一损失上的坐标下降求得；全部常数在四种语言端口中逐位一致。

相对 Mixbox~2.0，PBM 在深色加白套件上的平均 $\Delta E_{00}$ 为 2.258（最大 7.3），在留出深色加白对上为 2.377，在留出随机对上为 2.719--2.852，始终低于 RGB 插值与原始 K--M 叠加。相对独立的合成光谱世界，其各套件平均误差与参考相差至多 0.67 $\Delta E_{00}$，bootstrap 区间全部重叠；平均色相误差在六个套件中的三个更低。剩余的大误差（10.7--23.8 $\Delta E_{00}$）集中在"大变明度差叠加色相对立"的颜色对上，以及近中性深色盲区。

模型的实用价值在于它不需要什么：不需要光谱数据、不需要查找表、不需要训练的潜映射、也不需要第三方许可组件。因此它可以替换绘画软件、游戏引擎与网页应用中的线性 RGB 插值，同时复现颜料混色的定性行为——蓝加黄得到绿、互补色对得到浑浊中性色、深色加白仍保持色度。本文的评估把这一主张界定为"与参考以及与合成光谱世界一致"，并未确立与实测颜料一致。
